\chapter{Tools}
\label{ch:tools}

The Tools page groups the tools by what they work on: the \emph{device}, the
\emph{data}, and the \emph{test}. A dot next to a tool means it has something
going on (the simulator running, a scan in progress).

\screenshot[6cm]{tools}{The Tools page.}
  {The Tools page: the list of tools on the left (Device, Data, Test) and one
   tool open on the right.}

Tools that talk to a device use the connection of step 1 on the Main page, and
only read.

\section{Identify}

Asks the device for its vendor, product code and revision with \emph{Read
Device Identification} (FC43/14). Choose which \field{Objects} to ask for (basic,
regular or extended) and \uiaccent{Read identification}. Each object the device
returns is listed with its value.

\screenshot[4cm]{identify}{A device that answered FC43.}
  {The Identify tool after a read: vendor name, product code and revision.}

Not every device supports FC43: many answer exception 1 (illegal function).

\ui{Check function codes} sends one harmless request for each read function
(FC 1--4, 7, 8, 11, 17, 20, 24 and 43) and shows how the device answered:
\emph{Supported}; \emph{Not supported} when it answered ``illegal function'';
\emph{Supported, refused this request} when it answered another exception,
which means it has the function but not, for example, address 0; or
\emph{No answer}. Nothing is written. FC 7, 8, 11 and 17 exist on serial lines
only, so over TCP they often answer ``illegal function''.

The functions are tried at the \field{Address} given (0 at first). With
\ui{Also try write functions}, Modbench also tries FC 5, 6, 15, 16 and 23 by
reading coil and holding register \field{Address} and writing the same values
back, and FC 22 with a mask that keeps every bit. It always asks first, and the
option is off in read-only mode. A write whose value could not be read is
\emph{Not tried}. The values do not change, but a device can still act on a
write: only do this on a device you may change.

\ui{Find largest read} finds how many registers or bits the device returns in
one read from \field{Address}, by halving, up to 125 registers or 2000 bits,
and says why one more failed. ``Illegal data value'' means that is the
device's limit: set it as \field{Max registers} on the Main page. ``Illegal data
address'' means the area has a gap there.

\section{Find addresses}

Finds which addresses answer on a device without a manual. It reads blocks of
addresses, and splits the blocks that fail until it knows which addresses
exist.

\begin{description}[style=nextline, leftmargin=0pt, font=\intermedium]
  \item[Area, From, To] What to search.
  \item[Block] Addresses read at once before splitting. Larger blocks are
    faster on devices with long runs of registers.
  \item[Wait (ms)] How long to wait for each answer.
\end{description}

The result lists the ranges that answered. \ui{Read} next to a range sets up
a manual request for it on the Main page.

\screenshot[5cm]{find-addresses}{Address ranges found on an energy meter.}
  {Find addresses after a scan of holding registers 0--2000: the ranges that
   answer, with their counts.}

\section{Network scan}
\label{sec:network-scan}

Finds Modbus TCP devices on a subnet. It needs no connection on the Main page
and leaves an open one alone. Each address gets a TCP connection to the port;
when it opens, the scan asks for one holding register. Any answer, an exception
included, shows a Modbus device. Nothing is written.

\begin{description}[style=nextline, leftmargin=0pt, font=\intermedium]
  \item[Addresses] A subnet (\val{192.168.1.0/24}), a range
    (\val{192.168.1.10-50} or \val{192.168.1.10-192.168.2.20}), or addresses
    separated by commas. Up to 65536 addresses (a /16).
  \item[Port, Unit ID] Where to ask, and which unit ID the request carries.
  \item[Mode] \ui{Fast} waits a short time (150\,ms by default) and does not
    retry: a quick first pass to find the devices that are online.
    \ui{Thorough} waits longer (1000\,ms) and asks once more after a timeout,
    for slow devices, gateways and busy networks. The line under the settings
    says which wait and retries the scan will use.
  \item[Wait (ms)] How long to wait for the connection, and then for the
    answer, in the chosen mode; each mode keeps its own.
  \item[Send from] The network interface, shared with the connection of
    step 1.
  \item[Identify (FC43)] Also reads the vendor, product code and revision of
    the devices found.
\end{description}

Results appear as they are found. \ui{Modbus} marks a device that answered;
\ui{Port open} an address that accepted the connection but sent no Modbus
answer in time. \ui{Open} sets up the connection of step 1 for a device.
\ui{Stop} ends the scan and keeps what it found. An address that does not
answer within the wait is left out; that does not prove there is no device
there, so try a longer wait on a slow network.

\section{Exporting scan results}

The unit ID and serial settings scans (\ui{Find…} in step 2),
\ui{Find addresses} and \ui{Network scan} have \ui{Export CSV…} once there
are results. The file starts with a few lines that say which scan ran, when,
and with which settings and connection; a blank line; then one row per
result. A stopped scan says where it stopped. The unit ID scan lists every unit
ID, with \val{answered}, \val{no answer} or \val{not asked}; the other scans
list what answered and count the rest. The network scan can be exported while
it runs, without stopping it. Files are UTF-8, with the CSV separator of the
settings.

\section{Data decoder}
\label{sec:data-decoder}

Type or paste registers (in hex) to see them as every type at once, in all
four byte orders: signed and unsigned 16-, 32- and 64-bit integers, floats,
and text. \ui{Use last read} takes the registers last read on the Main page.

This is the quickest way to find a value's type and byte order before
writing a map: look for the column where the number makes sense (a voltage
near 230, a frequency near 50).

\screenshot[6cm]{data-decoder}{Two registers decoded in every type and order:
  only 32-bit float in ABCD order reads as a voltage, 230.4.}
  {The Data decoder with 4366 6666: every type in the four byte orders, and
   each register on its own.}

It works without a connection.

\section{Frame decoder}
\label{sec:frame-decoder}

Paste a frame from a log, a bus analyser or another tool (for example
\val{01 03 00 00 00 0A C5 CD}) to read it field by field, with its CRC or LRC
checked. Choose the \field{Framing} (RTU, ASCII, TCP) and who sent it
(\field{Client} or \field{Device}), since a request and an answer with the same
function code have different fields.

\screenshot[5cm]{frame-decoder}{An RTU request decoded.}
  {The Frame decoder with an RTU request: unit 1, function 3, address 0,
   quantity 10, CRC valid.}

\section{Trend player}
\label{sec:trend-player}

Plays back a CSV file made with \ui{Record…} on the Device map tab, so a test
can be looked at again later.

\begin{steps}
  \item \uiaccent{Open recording…} and choose the file.
  \item Tick the columns to plot (up to eight; text columns cannot be
        plotted).
  \item \ui{Play} at 1× to 60× speed, or drag the slider to any moment.
\end{steps}

\screenshot[7cm]{trend-player}{A recording played back.}
  {The Trend player with a one-hour recording: columns ticked, the slider
   half-way, speed 8×, and the charts.}

\field{Time span} works as on the Chart tab, and \ui{Export image…} saves the
charts (Section~\ref{sec:chart-image}).

\section{Statistics}
\label{sec:statistics}

Every request sent on the Main page since the last \ui{Reset}:

\begin{description}[style=nextline, leftmargin=0pt, font=\intermedium]
  \item[By result] How many succeeded, timed out, came back as each exception,
    or arrived as bad frames.
  \item[By function] Requests, failures, and minimum, average and maximum
    response time per function code.
\end{description}

A pattern points to the cause: timeouts to wiring, baud rate or unit ID;
exception 2 to an address the device does not have; bad frames to noise or
two devices with the same unit ID.

\screenshot[5cm]{statistics}{Statistics of a long test.}
  {The Statistics tool after a long test: results by cause, the table by
   function, and the suggested timeout with Use in step 1.}

\field{Suggested timeout} is four times the slowest answer;
\ui{Use in step 1} puts it in the connection settings.

\section{Alarms}

Limits on the registers of the open map. A value below \field{Low} or above
\field{High}, or a bit in its alarm state, turns red on the Device map tab.

\begin{steps}
  \item Open a map on the Main page.
  \item \ui{Add an alarm on a register…} and choose the register.
  \item Set the limits, or for a bit, the state that is an alarm.
\end{steps}

\screenshot[4cm]{alarms}{Alarms on a meter's map.}
  {The Alarms tool with three alarms: a voltage range and two upper limits,
   one of them exceeded and in red.}

\screenshot[5cm]{alarms-map}{The value in alarm, in red on the Device map tab.}
  {The Device map tab with an alarm: the current over its limit in red, and
   \emph{1 in alarm} above the table.}

Alarms are kept per map on this computer. Nothing is sent to the device.
